% ПРИЛОЖЕНИЕ. Изходният текст не се препечатва: дадени са модулите, ролята им и
% обемът им, плюс адресът на хранилището.
\section*{ПРИЛОЖЕНИЕ. ИЗХОДЕН ТЕКСТ}
\label{sec:appendix}
\addcontentsline{toc}{section}{ПРИЛОЖЕНИЕ}

Пълният изходен текст е в \url{https://github.com/Alex-Tsvetanov/waypoint} под лиценз MIT.
Табл.~\ref{tab:files} дава модулите, ролята им и обема им в редове, като заглавният файл и
файлът с реализация са събрани.

\begin{table}[!htb]
\caption{Модули на реализацията, 4776 реда общо}
\label{tab:files}
\centering
\fontsize{10}{12}\selectfont
\begin{spacing}{1.0}
\begin{tabular}{lL{7.6cm}c}
\toprule
Модул & Съдържание & Редове \\
\midrule
\code{types}    & скаларни типове, номер на възел, момент във времето & 38 \\
\code{env}      & абстракцията от пет операции над време и транспорт & 67 \\
\code{packet}   & формат по мрежата, кодиране и декодиране с проверка на границите & 376 \\
\code{lsdb}     & база с обявленията, поредни номера, възраст, оттегляне & 159 \\
\code{spf}      & граф и Дийкстра с многопътност при равна цена & 236 \\
\code{neighbor} & автомат за състоянията на съседите като чиста функция & 160 \\
\code{router}   & протоколно ядро: таймери, разпространение, маршрути & 621 \\
\code{topology} & четец на топологии и генератори за четирите семейства & 338 \\
\code{sim}      & дискретно-събитиен симулатор, модел на връзките, откази & 471 \\
\code{analysis} & сходимост, преходни цикли, натоварване, нестабилност & 104 \\
\code{udp}      & сокети и таймери в жив режим, Winsock и BSD & 351 \\
\code{apps/demo}  & сценарият с времевата картина и изход в DOT & 239 \\
\code{apps/bench} & четирите експеримента и записът в CSV & 326 \\
\code{apps/live}  & демонът, един процес на маршрутизатор & 153 \\
\code{tests/}     & собствен изпълнител на тестове и осем набора, 70 случая & 1137 \\
\bottomrule
\end{tabular}
\end{spacing}
\end{table}

Всичките 250 резултата са валидни, така че таблиците в раздел~\ref{sec:results} съдържат пълния
набор. Измерванията се възпроизвеждат с \code{cmake -S . -B build}, \code{cmake -{}-build build}
и \code{./build/waypoint-bench -{}-csv results.csv}.
